% =========================================================
% Conclusion — short, declarative; mirrors benchmark framing
% =========================================================
\section{Conclusion}
We introduce \textsc{MineNPC-Task}, a practical benchmark for evaluating mixed-initiative, memory-aware LLM agents in \emph{Minecraft} using only public in-game interfaces. 

Tasks are elicited from expert co-play, normalized into compact templates with explicit preconditions, and paired with simple validators that judge completion from bounded, in-world evidence. 

An initial snapshot with GPT-4o over 44 tasks (216 subtasks) highlights where agents struggle slot clarification, preconditions, memory reuse, and code execution yielding a \(\sim\)33\% subtask failure rate. 

We hope this benchmark encourages models and methods that plan across dependent steps, ask when uncertain, and ground memory in observable state without relying on hidden shortcuts.

% =========================================================
% Limitations — crisp, scoped; avoids over-claiming
% =========================================================
\section{Limitations}

\textbf{Scope.} Results apply to \emph{Minecraft} with a Mineflayer client under a bounded-knowledge policy. Portability to other engines or sensing/action envelopes is not evaluated.

\textbf{Model coverage.} We report a single-model snapshot (GPT-4o). There are no ablations or cross-model comparisons, so relative performance remains an open question.

\textbf{Task coverage.} The suite is expert-elicited and modest in size. It captures common goals we observed but does not exhaust the space of open-world play; selection bias is possible.

\textbf{Measurement granularity.} Validators return pass/fail from in-world traces. This improves reproducibility but under-represents partial progress and does not yet score cost or efficiency.

\textbf{Run-time variance.} Live co-play introduces practical variability (e.g., pathing, chunk loading). Policies reduce but do not eliminate noise; we report denominators to contextualize outcomes.

\textbf{Interaction design choices.} The framework prioritizes a single clarifying question and short plan previews. Richer multi-turn clarification or alternative planning styles may yield different results.

\textbf{Telemetry and auditability.} Released logs emphasize end-state evidence and high-level traces. Finer-grained diagnostics (e.g., pre-execution static checks) are limited and left to future iterations.

\textbf{Generalization of corrections.} While preferences and names are stored, persistent generalization of user corrections across sessions is limited in the current release.

% =========================================================
% Future Work — positive scope; actionable and aligned with benchmark goals
% =========================================================
\section{Future Work}

\textbf{Comparable baselines.} Add multi-model evaluations under identical perception–action contracts and policies, enabling controlled comparisons and public leaderboards.

\textbf{Expanded task pool.} Grow the expert-derived templates with parameterized variants and richer dependencies (e.g., multi-session builds), while keeping validators simple and reproducible.

\textbf{Richer metrics.} Report partial credit and efficiency (time, distance, resource cost), number of clarifications, and repair rate to complement pass/fail outcomes.

\textbf{Pre-execution checks and telemetry.} Introduce lightweight static checks for common API/parameter faults and expose finer-grained traces to improve diagnosis and replication.

\textbf{Targeted robustness probes.} Add focused tests for egocentric references and tool–affordance errors; surface memory contents and staleness in the UI to align expectations.

\textbf{Artifacts and reproducibility.} Continue releasing templates, validators, prompts, and logs so others can rerun the suite, contribute tasks, and extend the benchmark over time.
